\section{A same-access separation between a decision and an absolute forecast}
\label{sec:sameaccesslower}

This lower bound concerns acquired information, not the running time of a
solver given exact weights.  It grants the forecaster strictly more
observations than the finite-hold decision needs.  In particular it grants
perfect clocks, arbitrary calibrated additive fields, and complete continuous
observation of each released local cell.  The obstruction therefore does not
come from a slow unknown clock or from requiring an inefficient parameter
estimator.

\paragraph{The experiment class.}
There are $n$ independent binary components, with unknown logits
$\theta_i$, equilibrium failure probabilities $q_i=\sigma(\theta_i)$,
and \emph{known unit refresh clocks}.  At applied field $u$, a local cell
has rates
\[
  0\longrightarrow1:\ \sigma(\theta_i+u),\qquad
  1\longrightarrow0:\ 1-\sigma(\theta_i+u).
\]
An experiment may adaptively choose a component, an arbitrary field, a
known initial bit, a release duration, and its stopping rule.  It may observe
the complete resulting local path.  Several cells may operate in parallel;
their free durations are added.  A supplied fresh equilibrium bit at a
chosen field is also permitted and counted as one local observation.  Pins,
controller randomization, and computations convey no additional unknown
parameter by themselves.  An equilibrium bank of $n$ bits counts as $n$
local observations; parallelism does not turn its $n$ sensing channels into
one unit of acquired evidence.  These are the declared local state/clock
experiments, not direct measurements of unknown Hamiltonian coefficients.

Write
\[
  W=B+r^*\sum_{\text{released cells}} \text{free duration},
\]
where $B$ is the number of supplied fresh equilibrium bits and $r^*$ is
the known maximum available refresh clock (one in the unit-clock model).
Allowing additional known clock modules up to a fixed $r^*$ only changes
this normalization: their rates are the same under the two hypotheses.
This relaxed
bill omits pin, detector, field, memory, and computational costs, so a lower
bound on it remains a lower bound when those services are charged.

\paragraph{Theorem (forecast information lower bound).}
Let $n\ge4$, $0<\epsilon\le1/8$, and $0<\delta<1/2$.
Suppose an adaptive procedure returns an estimate $\widehat M$ satisfying
\[
 \Pr_\theta\{|\log\widehat M-\log M(\theta)|\le\epsilon\}\ge1-\delta
\]
for every allowed model, where $M$ is the mean first visit to
$b=(1,\ldots,1)$ from the equilibrium law conditioned on $b^c$.
Then, on at least one allowed model,
\begin{equation}
  \mathbb E W\ \ge
  \frac{n^2}{512e\,\epsilon^2}
       \operatorname{kl}(1-\delta,\delta).
  \label{eq:forecastlower}
\end{equation}
With only fresh equilibrium-bit experiments the stronger constant
$n^2\operatorname{kl}(1-\delta,\delta)/(128\epsilon^2)$ holds.
For fixed $\delta<1/2$, this is an $\Omega(n^2/\epsilon^2)$ local-evidence
requirement.  It holds even though the operating gap equals one, the clocks
are known, and every coordinate can be actively balanced.

\paragraph{Proof.}
First consider equal logits $\theta_i=\theta$, and put $q=\sigma(\theta)$,
$p=q^n$.  Independence and the two-state semigroup give
\[
 P_t(b,b)=\{q+(1-q)e^{-t}\}^n.
\]
The classical fundamental-matrix identity, also derived in the return
certificate below, gives
\begin{align}
 S_n(q)&=\int_0^\infty[P_t(b,b)-p]\,dt
       =\sum_{k=1}^n {n\choose k}
                    \frac{q^{n-k}(1-q)^k}{k},\\
 M_n(q)&=\frac{S_n(q)}{p(1-p)}.
 \label{eq:productmeanlower}
\end{align}
Let $K\sim\operatorname{Bin}(n,1-q)$.  Then
$S_n=\mathbb E[K^{-1};K>0]$.  Since
$1/k\le2/(k+1)$ for $k\ge1$,
\[
 S_n(q)\le\frac{2}{(n+1)(1-q)}.
\]
Direct differentiation of the finite sum gives
\begin{equation}
 -\frac{d}{d\theta}\log M_n(\sigma(\theta))
 =\frac{1-p}{S_n(q)}-\frac{np(1-q)}{1-p}.
 \label{eq:meanlogderivative}
\end{equation}
For $n\ge4$ and $q\le1/2$, the right side is at least
\[
 \frac{15(n+1)}{64}-\frac n{15}\ \ge\ \frac n8.
\]
Take $\theta^{(0)}_i=-\log3$ and
$\theta^{(1)}_i=-\log3+d$, where $d=32\epsilon/n\le1$.
The probabilities along this segment stay below $1/2$.
Equation~\eqref{eq:meanlogderivative} therefore yields
\[
 \log M(\theta^{(0)})-\log M(\theta^{(1)})\ge4\epsilon.
\]
A valid estimate distinguishes these two hypotheses with both error
probabilities at most $\delta$.

For one equilibrium-bit observation at any adaptively selected field,
the Bernoulli log-partition function has second derivative at most $1/4$.
Consequently its conditional divergence between the two hypotheses is at
most $d^2/8$, independently of the field magnitude.  For a released cell,
write $\lambda_0(x,u),\lambda_1(x,u)$ for the applicable exit rates.
Because the derivative of each logistic log rate has magnitude at most one,
\[
 \left|\log\frac{\lambda_1(x,u)}{\lambda_0(x,u)}\right|\le d.
\]
The continuous-path likelihood identity gives conditional divergence
\[
 \mathbb E_0\int
 \left[\lambda_0\log\frac{\lambda_0}{\lambda_1}
             -\lambda_0+\lambda_1\right]dt
 \le\frac{ed^2}{2}\,\mathbb E_0\int dt,
\]
since $\lambda_0\le1$ in the unit-clock case (or at most $r^*$ with
additional clock modules) and
$e^v-1-v\le e d^2/2$ for $|v|\le d\le1$.
Adaptive fields, adaptive stopping, and repeated use of a cell obey the
same conditional likelihood chain rule.  Truncating an unbounded stopping
rule and then taking a limit gives
\[
 \operatorname{KL}(\mathsf P_0^{\rm data}\Vert\mathsf P_1^{\rm data})
 \le\frac{ed^2}{2}\,\mathbb E_0 W.
\]
Data processing through the final binary test lower-bounds this divergence
by $\operatorname{kl}(1-\delta,\delta)$, proving
\eqref{eq:forecastlower}.  The equilibrium-only bound uses $d^2/8$ in
place of $ed^2/2$.  Finite instrument errors cannot increase the information
available beyond this deliberately idealized observation class.
\hfill$\square$

\paragraph{Heterogeneous and nontrivial decision instances.}
The hard pair need not consist of identical components.  Let initial
$q_i$ belong, for example, to $\{1/4,1/3\}$ in any unknown arrangement.
Apply the same additive logit shift $t$ to every component.
If $K$ is the resulting Poisson-binomial number of working components,
$S=\mathbb E[K^{-1};K>0]$ and $p=\prod_iq_i$ satisfy
\[
 \frac{dS}{dt}=\Lambda S-(1-p),\qquad
 \frac{d\log p}{dt}=\Lambda,
 \qquad \Lambda=\sum_i(1-q_i).
\]
Thus
\[
 -\frac{d\log M}{dt}=\frac{1-p}{S}-\frac{p\Lambda}{1-p}.
\]
As long as all $q_i\le1/2$, stochastic domination by
$\operatorname{Bin}(n,1/2)$ gives the same lower bound $n/8$.
For $n\ge8$, $\epsilon\le1/8$, the shift $32\epsilon/n$ keeps the
displayed heterogeneous family below $1/2$.  The information bound is
unchanged.  These models can have a nontrivial optimal assignment of unequal
clock modules, while the unknown common energy shift leaves their optimal
ordering unchanged.  The lower bound is therefore not an artifact of
uncertainty about a globally rescaled time unit.

\paragraph{Matching order for prediction in this product family.}
The $n^2$ dependence is attainable up to confidence and finite-service
logarithms; it is not an inverse-$p$ assertion.  For $q_i$ in a fixed compact
subinterval of $(0,1)$, sample each local equilibrium probability $m$ times,
clip its empirical frequency to a slightly larger compact subinterval,
and add the $n$ log frequencies.  The bias is $O(n/m)$ and bounded
differences give fluctuation
$O(\sqrt{n\log(1/\delta)/m})$.  Taking
$m=O(n\epsilon^{-2}\log(1/\delta))$ estimates $\log p$ to $O(\epsilon)$
using $O(n^2\epsilon^{-2}\log(1/\delta))$ local observations.  The
nonrare return prefactor $S$ is acquired separately by the finite return
operation.  Every ideal equilibrium read can be realized by a finite hold
at the known unit clock, with conditional bias below its allotted error;
this introduces an explicit logarithmic dwell factor, not $1/p$.

\paragraph{Meaning of the separation.}
At fixed margins and fixed decision accuracy, the direct acquired
intervention constructions use $\widetilde O(n)$ local evidence, while
\emph{every} uniformly accurate absolute-mean forecaster in the same local
access class has the $\Omega(n^2)$ requirement above.  A method is free to
use the same nuisance cancellation and become a direct decision method;
the theorem does not claim a lower bound against such a method.  It shows
that requiring a calibrated absolute forecast before deciding is an
additional statistical task.  It is not a claim that ordinary simulation
is the strongest baseline, nor that all rare-event decisions are cheap.

The proof uses classical likelihood/change-of-measure tools, as in
fixed-confidence experimental identification.  The established
fundamental-matrix and bandit results are machinery; no priority claim is
made for those identities.  The derived separation is for the explicit
rare first-event intervention family and the paid local experiment class
above.
